\chapter{Cuando algo falla, y dónde se despliega}
\label{ch:resiliencia}

\section{El hub desaparece}

El cliente decide con estado local, así que un hub caído no detiene el tráfico. Lo que cambia con el tiempo es
cuánto se fía el cliente de lo último que supo.

Se para el hub, y solo el hub:

\begin{lstlisting}[style=sh]
kill $(lsof -nP -iTCP:18443 -sTCP:LISTEN -t)
\end{lstlisting}

\subsection{Primera fase: no pasa nada}

El último listado aceptado vale hasta el 13 de noviembre a las 00:00:05. Doce horas antes:

\begin{lstlisting}[style=sh]
clock 2026-11-12T12:00:00Z; sleep 2
curl -s $ADMINAPI/status | python3 -c "import json,sys;d=json.load(sys.stdin);print(d['state'], d['last_accepted_listing'])"
row m1-forecast
\end{lstlisting}
\begin{lstlisting}[style=out]
Active {'sequence': 15, 'valid_from': '2026-11-12T00:00:05Z', 'valid_until': '2026-11-13T00:00:05Z'}
m1-forecast   HTTP/1.1 200 OK | -
\end{lstlisting}

El estado sigue siendo \cod{Active}. El fallo de sincronización queda registrado, pero no cambia nada:

\begin{lstlisting}[style=json]
{"type":"vhub.sync","ts":"2026-11-12T12:00:00Z","code":"SYNC_UNREACHABLE","severity":"error",
 "detail":"Get \"https://127.0.0.1:18443/v1/namespaces/flexgrid/versions\": dial tcp 127.0.0.1:18443: connect: connection refused"}
\end{lstlisting}

\subsection{Segunda fase: \texttt{Degraded}}

Pasada la caducidad del listado:

\begin{lstlisting}[style=sh]
clock 2026-11-13T06:00:00Z; sleep 2
row m1-forecast -b
\end{lstlisting}
\begin{lstlisting}[style=out]
m1-forecast   HTTP/1.1 400 Bad Request | VHub-Verdict: violation;gate=policy_intake;code="CLIENT_DEGRADED_REJECT";rule="contract-request"
    {"@context":["https://w3id.org/dspace/2025/1/context.jsonld"],"@type":"ContractNegotiationError",
     "code":"CLIENT_DEGRADED_REJECT", … ,"reason":["CLIENT_DEGRADED_REJECT"]}
\end{lstlisting}

El estado es \cod{Degraded} y \cod{/readyz} sigue respondiendo \cod{200}. Un mensaje que ayer era válido hoy se
rechaza, y el código explica que no es culpa del mensaje.

Es exactamente lo que el \cod{SpaceProfile} pidió en el capítulo~\ref{ch:spaceprofile}:

\begin{lstlisting}[style=yaml]
  degraded:
    policy_intake: reject
    dsp_message: continue
\end{lstlisting}

Sin contacto con el hub, Flexgrid deja de admitir \textbf{políticas nuevas}, porque podría haberse revocado
algo sin que el cliente lo sepa. Pero un gate con \cod{continue} sigue clasificando con normalidad. Una
transferencia mal formada lo demuestra, porque se sigue evaluando contra el esquema de DSP:

\begin{lstlisting}[style=out]
t2-transfer-no-format   HTTP/1.1 400 Bad Request | VHub-Verdict: violation;gate=dsp_message;code="VERDICT_VIOLATION";rule="transfer-request", violation;gate=agreement_admission;code="AGREEMENT_VERSION_REVOKED";rule="transfer-request"
\end{lstlisting}

La cabecera lleva aquí dos elementos, uno por gate. Cuando una regla tiene varios gates se toma la acción más
restrictiva, y cada gate que no es \ok\ aporta su veredicto.

\begin{ideaClave}
\cod{degraded} es una decisión de negocio que toma la autoridad, no cada participante: viaja en el manifiesto
firmado. Responde a la pregunta «si no sé si las reglas han cambiado, ¿qué es peor: parar o seguir?», y la
respuesta puede ser distinta para cada gate.
\end{ideaClave}

\subsection{Tercera fase: \texttt{FailClosed}}

Pasado \cod{max_staleness}, siete días, desde la caducidad:

\begin{lstlisting}[style=sh]
clock 2026-11-21T06:00:00Z; sleep 2
row m1-forecast
curl -s -o /dev/null -w '%{http_code}\n' $ADMINAPI/readyz
\end{lstlisting}
\begin{lstlisting}[style=out]
m1-forecast   HTTP/1.1 400 Bad Request | VHub-Verdict: violation;gate=policy_intake;code="CLIENT_FAIL_CLOSED";rule="contract-request"
503
\end{lstlisting}

Ahora todos los gates rechazan, digan lo que digan sus ajustes, y \cod{/readyz} pasa a \cod{503}. Los cambios de
estado quedan como eventos:

\begin{lstlisting}[style=json]
{"type":"vhub.state","ts":"2026-11-13T06:00:00Z","severity":"info","detail":"state changed","from":"Active","to":"Degraded"}
{"type":"vhub.state","ts":"2026-11-21T06:00:00Z","severity":"info","detail":"state changed","from":"Degraded","to":"FailClosed"}
\end{lstlisting}

\subsection{La recuperación}

\begin{lstlisting}[style=sh]
$BIN/alexandria-server -config server.yaml -clock-file clock >> server.log 2>&1 &
clock 2026-11-21T06:10:00Z; sleep 3
curl -s $ADMINAPI/status | python3 -c "import json,sys;d=json.load(sys.stdin);print(d['state'], d['last_accepted_listing'])"
row m1-forecast
\end{lstlisting}
\begin{lstlisting}[style=out]
Active {'sequence': 16, 'valid_from': '2026-11-21T06:00:00Z', 'valid_until': '2026-11-22T06:00:00Z'}
m1-forecast   HTTP/1.1 200 OK | -
\end{lstlisting}

Al volver, el hub emite un listado nuevo, de secuencia 16. El cliente reintenta un minuto después del último
fallo, lo acepta y vuelve a \cod{Active} sin intervención. No hay que reiniciar nada.

\begin{trampaHabitual}[Con reloj de fichero, el reintento también espera al reloj]
El reintento se cuenta en tiempo del reloj, no de pared. Si arrancas el hub y no mueves el fichero, el cliente
no reintenta nunca: para él no ha pasado ni un segundo. Por eso el comando anterior avanza diez minutos.
\end{trampaHabitual}

\begin{center}
\begin{tikzpicture}[font=\scriptsize, x=1.45cm, y=1cm]
  \draw[alxBorder, -{Stealth[length=2mm]}] (0,0) -- (9.6,0);
  \fill[alxTeal!35] (0,0.15) rectangle (3,0.75);    \node at (1.5,0.45) {\texttt{Active}};
  \fill[alxAmber!40] (3,0.15) rectangle (7,0.75);   \node at (5,0.45) {\texttt{Degraded}};
  \fill[alxRed!30] (7,0.15) rectangle (8.6,0.75);   \node at (7.8,0.45) {\texttt{FailClosed}};
  \fill[alxTeal!35] (8.6,0.15) rectangle (9.5,0.75);
  \draw (1,0.1) -- (1,-0.1) node[below] {hub cae};
  \draw (3,0.1) -- (3,-0.1) node[below] {\texttt{validUntil}};
  \draw (7,0.1) -- (7,-0.1) node[below] {+ \texttt{max\_staleness}};
  \draw (8.6,0.1) -- (8.6,-0.1) node[below] {listado nuevo};
  \node[above] at (5,0.78) {por gate: \texttt{continue} o \texttt{reject}};
  \node[above] at (7.8,0.78) {todo se rechaza};
\end{tikzpicture}
\end{center}

\section{Cuando el hub responde, pero mal}

Un hub que no contesta es un problema de disponibilidad. Uno que contesta algo incorrecto es un problema de
seguridad, y el cliente lo trata con más dureza: descarta el listado entero y sigue con la instantánea anterior.

\noindent\begin{tabularx}{\linewidth}{@{}L{6.7cm}Y@{}}
\toprule
\textbf{Código} & \textbf{Qué ha visto el cliente} \\
\midrule
\cod{SYNC_SIGNATURE_INVALID} & La firma del listado no verifica contra el DID del hub. \\
\cod{SYNC_ROLLBACK} & Una secuencia menor que la última aceptada: alguien sirve un listado antiguo. \\
\cod{SYNC_SEQUENCE_CONFLICT} & La misma secuencia con otro contenido. \\
\cod{SYNC_REVOCATION_DROPPED} & Una revocación que estaba ha desaparecido. \\
\cod{SYNC_DATES_CHANGED} & Una fecha ya emitida ha cambiado. \\
\cod{SYNC_MANIFEST_DIGEST_MISMATCH}, \cod{SYNC_FILE_HASH_MISMATCH} & Un paquete o un fichero no coincide con su resumen. \\
\bottomrule
\end{tabularx}

\medskip
Estos casos no se han provocado en este tutorial, porque exigen un hub que mienta. La batería de conformidad
del repositorio tiene un caso negativo mínimo para cada uno:

\begin{lstlisting}[style=sh]
go test ./conformance -run 'TestN_' -v
\end{lstlisting}

\section{Mensajes problemáticos}

\noindent\begin{tabularx}{\linewidth}{@{}L{5.5cm}L{5.9cm}C{2.2cm}@{}}
\toprule
\textbf{Mensaje} & \textbf{Código, con la regla en \texttt{enforce}} & \textbf{¿Ejecutado?} \\
\midrule
Política fuera del esquema & \cod{VERDICT_VIOLATION} & \si \\
Contexto JSON-LD que no está en el paquete & \cod{CONTEXT_NOT_IN_PACKAGE} & \si \\
Envoltura DSP incompleta & \cod{VERDICT_VIOLATION} en \cod{dsp_message} & \si \\
Ninguna versión vigente todavía & \cod{NO_VERSION_IN_FORCE} & \si \\
Acuerdo de una versión revocada & \cod{AGREEMENT_VERSION_REVOKED} & \si \\
Listado caducado, gate en \cod{reject} & \cod{CLIENT_DEGRADED_REJECT} & \si \\
Listado caducado más de \cod{max_staleness} & \cod{CLIENT_FAIL_CLOSED} & \si \\
Cuerpo que no es JSON, o sin el nodo de \cod{select} & \cod{MESSAGE_UNPARSEABLE} & no \\
Petición sin regla, con \texttt{default\_action: reject} & \cod{RULE_DEFAULT_REJECT} & no \\
\bottomrule
\end{tabularx}

\section{Dos formas de desplegar el cliente}

El cliente decide sobre cada mensaje. \emph{Quién reenvía} el mensaje es una decisión de despliegue, que se
elige con \cod{mode}.

\begin{center}
\begin{tikzpicture}[font=\small, node distance=7mm and 13mm,
  box/.style={draw=alxNavy, fill=alxLightBlue, rounded corners=2pt, minimum height=9mm, minimum width=29mm, align=center},
  ext/.style={draw=alxBorder, fill=alxSlate, rounded corners=2pt, minimum height=9mm, minimum width=24mm, align=center},
  arr/.style={-{Stealth[length=2.2mm]}, alxNavy, thick}]
  \node[ext] (p1) {Contraparte};
  \node[box, right=of p1] (c1) {\texttt{alexandria-client}};
  \node[ext, right=of c1] (k1) {Conector};
  \draw[arr] (p1) -- (c1); \draw[arr] (c1) -- (k1);
  \node[left=3mm of p1, anchor=east] {\texttt{proxy}};

  \node[ext, below=17mm of p1] (p2) {Contraparte};
  \node[box, right=of p2, fill=alxSlate, draw=alxBorder] (e2) {Envoy};
  \node[ext, right=of e2] (k2) {Conector};
  \node[box, below=8mm of e2] (c2) {\texttt{alexandria-client}};
  \draw[arr] (p2) -- (e2); \draw[arr] (e2) -- (k2);
  \draw[arr] (e2) to[bend right=18] node[left]{\scriptsize\texttt{ext\_authz}} (c2);
  \draw[arr] (c2) to[bend right=18] node[right]{\scriptsize 200 o rechazo} (e2);
  \node[left=3mm of p2, anchor=east] {\texttt{validator}};
\end{tikzpicture}
\end{center}

\noindent\begin{tabularx}{\linewidth}{@{}L{3.3cm}Y Y@{}}
\toprule
 & \textbf{\texttt{mode: proxy}}, por defecto & \textbf{\texttt{mode: validator}} \\
\midrule
Quién reenvía & El cliente. & Un proxy externo, como Envoy. \\
Papel del cliente & Decide y reenvía. & Solo decide; nunca toca el conector. \\
\cod{upstreams} & Obligatorio. & No se usa. \\
\emph{Listeners} de tráfico & Proxies inverso y directo. & Puntos de consulta a los que llama el proxy. \\
\bottomrule
\end{tabularx}

\medskip
El camino de decisión es el mismo código en los dos modos: reglas, gates, clasificación, cabecera, eventos,
métricas y API de administración. En modo \cod{validator}, el cliente responde \cod{200} para permitir y el
rechazo habitual, con cuerpo y cabecera, para denegar. Es el contrato del filtro \cod{ext_authz} de Envoy.

\begin{lstlisting}[style=yaml]
mode: validator
trust:
  hub: "https://127.0.0.1:8443"
  namespace: "demo"
  hub_did: "did:key:z6Mk…"
listeners:
  dsp-ingress: { bind: "127.0.0.1:9085" }   # solo debe alcanzarlo el proxy
  admin: { bind: "127.0.0.1:9090" }
rules:
  seed: "file:deploy/rules.seed.json"
events:
  sink: "file:var/client/events.jsonl"
\end{lstlisting}

Lo que cambia en modo \cod{validator}, porque el cliente ya no ve la respuesta del conector:

\begin{itemize}[leftmargin=*, itemsep=2pt]
  \item La \textbf{captura de acuerdos} ocurre cuando el mensaje se permite, no cuando el conector responde
        \cod{2xx}.
  \item Las reglas con \cod{validate_response} no se evalúan.
  \item El \textbf{reenvío byte a byte} pasa a ser responsabilidad del proxy.
  \item La \textbf{disponibilidad} es un ajuste del proxy: el ejemplo de Envoy cierra el paso si el validador
        no responde.
\end{itemize}

El repositorio trae un ejemplo ejecutable de cada topología, que levanta su propio laboratorio y repite el
escenario F6 comprobando la misma tabla de 29 veredictos:

\begin{lstlisting}[style=sh]
go run ./examples/proxy -bin-dir bin
\end{lstlisting}
\begin{lstlisting}[style=out]
…
2026-11-12T00:00:10Z F6.4 A1 transfer    400 AGREEMENT_VERSION_REVOKED    violation;gate=agreement_admission;code="AGREEMENT_VERSION_REVOKED";rule="transfer-request"
F6 scenario completed: 29/29 rows verified.
\end{lstlisting}

\begin{notaPractica}[Qué no se ha ejecutado]
\texttt{go run ./\allowbreak{}examples/\allowbreak{}envoy} necesita Envoy en el \cod{PATH} o descargar su imagen de Docker, y no se ha
ejecutado para este documento. Lo que se dice aquí del modo \cod{validator} procede de \cod{examples/README.md},
de \cod{examples/envoy/README.md} y del código de configuración.
\end{notaPractica}

\section{Cerrar el laboratorio}

\begin{lstlisting}[style=sh]
for port in 18443 19400 18282 19090; do kill $(lsof -nP -iTCP:$port -sTCP:LISTEN -t) 2>/dev/null; done
rm -rf var clock *.token *.log last-candidate.* client.local.yaml idp.key
\end{lstlisting}

Borrar \cod{var/} elimina la base de datos, las claves y el estado del cliente. El siguiente arranque genera un
DID nuevo.

\begin{retos}
\begin{enumerate}[leftmargin=*, itemsep=2pt]
  \item Cambia \cod{degraded.policy_intake} a \cod{continue}, publícalo como una versión nueva y repite la caída
        del hub. ¿Qué cambia en la segunda fase? ¿Y en la tercera?
  \item Con el hub caído, intenta cambiar una regla del cliente. ¿Depende del hub?
  \item Ejecuta \texttt{go run ./\allowbreak{}examples/\allowbreak{}proxy -keep} y envía \cod{m1-forecast} a su \emph{ingress}. ¿Por qué se
        rechaza con ese laboratorio y no con el tuyo?
\end{enumerate}
\end{retos}
